\documentclass[11pt,a4paper]{article}
% \documentclass{elsarticle}

\usepackage{amsmath,amssymb,amsfonts,graphicx,hyperref}
\usepackage[margin=1in]{geometry}


\usepackage{amstext}
\usepackage{stmaryrd}
%\usepackage{listings}
%\lstset{language=C++}
%\usepackage[boxed]{algorithm2e}
\usepackage{algpseudocode}

\usepackage[normalem]{ulem}
\usepackage{graphicx}

%\usepackage{authblk}
%\usepackage{xcolor}
\usepackage{array}
\usepackage{verbatim}
\usepackage{booktabs}
\usepackage{calc}
\usepackage{textcomp}
\usepackage{multirow}
%\usepackage[boxed]{algorithm2e}

\usepackage[all]{xy}

\usepackage{cancel}
\usepackage{mathtools}
\usepackage{placeins}
\usepackage{xspace}

%\input{packages}

\usepackage{tikz,tikzscale}
\usetikzlibrary{calc}
\usetikzlibrary{shadings}
%\usepackage[backend=biber,style=alphabetic]{biblatex}
%\addbibresource{literature.bib}

%%%%%%%%%% Packages to help editing %%%%%%%%%%
%\usepackage[notref,notcite]{showkeys} % Show labels.
\usepackage[mode=multiuser,status=draft]{fixme} % Provide note, warning, error, fatal.
%\fxsetup{layout=inline}%pdfcsignote}
\fxsetup{envlayout=color}
\fxsetup{innerlayout=inline}
\fxsetup{targetlayout=colorcb}
\FXProvidesTargetLayout{color}

\fxusetheme{color}
\FXRegisterAuthor{aa}{envmr}{AA}
\FXRegisterAuthor{mw}{envmw}{MW}

\usepackage{subcaption}
\usepackage[format=plain,indention=.5cm]{caption}
\usepackage{pgfplots}
\pgfplotsset{compat=1.18}
\usepgfplotslibrary{fillbetween}

\begin{comment}

\DontPrintSemicolon
\SetKwComment{Comment}{$\triangleleft$ }{}
\SetKwProg{Fn}{Function}{}{end}

\providecommand{\keywords}[1]
{
    \small
    \textbf{\textit{Keywords:}} #1
}
\end{comment}
% \newcommand{\FigureSubst}[1]{ \fbox{ 
% \begin{minipage} {0.7\textwidth} #1  
% \end{minipage}  }   }

\def\mesh{\mathbb M}
\def\cell{K}
\def\R{\mathbb R}
\def\ncol{N_{\text{colors}}}
\def\rest#1{I_{#1}^\downarrow}
\def\prol#1{I_{#1}^\uparrow}

\usepackage{colortbl} % Define a row color
\newcommand{\myrowcolor}{\rowcolor[gray]{0.925}}
\newcommand{\headercolor}{\rowcolor[gray]{0.825}}
\usepackage{booktabs} % Allows the use of \toprule, \midrule and \bottomrule in tables   

\newcommand{\pluseq}{\stackrel{+}{\gets}}
%\newcommand{\pluseq}{\stackbin{+}{=}}

\newcommand{\Dd}{\text{d}}

% Mohsen's setting and definitions
%###################################%
\usepackage{bm,upgreek}
\newcommand\rd{{\rm{d}}}
\newcommand{\bfm}[1]{{\rm\bf #1}}

\newcommand{\greentext}[1]{{\color{darkgreen} #1}}
\newcommand{\bluetext}[1]{{\color{blue} #1}}
\newcommand{\redtext}[1]{{\color{red} #1}}

\DeclareMathOperator{\tr}{tr}
%###################################%

% MW settings
%###################################%

\newcommand{\DG}{{\mathbb{V}_h}}
\newcommand{\DGell}{{\mathbb{V}_\ell}}
%###################################%
%

\usepackage{graphicx} % Required for inserting images
\usepackage[linesnumbered,vlined,commentsnumbered,ruled]{algorithm2e}

% \newcommand{\revOld}[1]{}
% \newcommand{\revNew}[1]{{{#1}}}
\newcommand{\revNew}[1]{{\color{blue}{#1}}}   % only in original submission
\newcommand{\revOld}[1]{\sout{\protect\color{red}{#1}}}  % only in rev 1

\definecolor{darkgreen}{rgb}{0.0, 0.5, 0.0}

% \usepackage[marginal]{showlabels}

\definecolor{deepgreen}{RGB}{0,100,0}

% --- plot size / style helpers (added) ---
% Define global plot sizes (adjust here to change all figures)
\newcommand{\plotwidth}{0.49\columnwidth}
\newcommand{\plotheight}{0.4\columnwidth}

% Create a cycle list that uses symbols (markers) for the paper plots
\pgfplotscreateplotcyclelist{paper markers}{
    {teal,    solid, mark=*},
    {orange,   solid, mark=triangle*},
    {blue,     solid, mark=square*},
    {red,      solid, mark=pentagon*},
    {violet,     solid, mark=otimes*},
    {brown,    solid, mark=halfcircle*},
    {black,   solid, mark=diamond*},
}

% Default to the solid paper markers for all plots.
\pgfplotsset{cycle list name=paper markers}

% Centralized palette (change these three names to update all highlights)
\newcommand{\PaperColorA}{teal}
\newcommand{\PaperColorB}{orange}
\newcommand{\PaperColorC}{blue}

% Global opacity for range highlights (change here)
\newcommand{\RangeFillOpacity}{0.15}

% Reusable style for axis environments in the paper
\pgfplotsset{
    paperplot/.style={
            width=\plotwidth,
            height=\plotheight,
            cycle list name=paper markers,
            every axis/.append style={font=\small},
            grid=major,
            thick
        }
}

\pgfplotsset{
    longplot/.style={
            width=0.8\columnwidth,
            height=\plotheight,
            cycle list name=paper markers,
            every axis/.append style={font=\small},
            grid=major,
            thick
        }
}


\title{A Geometric Multigrid Preconditioner for Shifted Boundary Method}
\author{Micha\l{} Wichrowski\footnote{Heidelberg University, Germany, mt.wichrowsk@uw.edu.pl}, Ajay Ajith
    \footnote{Heidelberg University, Germany}}
% \\date{\\today}

\begin{document}

% \begin{abstract}

% \end{abstract}
% \begin{keyword}
%     immersed boundary methods \sep finite element methods \sep Shifted Boundary Method \sep Discontinuous Galerkin \sep
%     multigrid
% \end{keyword}
\maketitle
\section{Introduction}
\mwerror{Introduction needs to be rewritten.}

\section{Method formulation}
\label{sec:method}
We consider the Poisson problem as a model problem:
\begin{align}
    -\Delta u & = f \quad \text{in } \Omega, \label{eq:poisson}      \\
    u         & = g \quad \text{on } \Gamma, \label{eq:dirichlet_bc}
\end{align}
where $\Omega \subset \mathbb{R}^d$ ($d=2,3$) is a domain with boundary $\Gamma =\partial\Omega$ as depicted in
Figure~\ref{fig:sbm_setup}, $f$ is a given source term, and $g$ is the prescribed Dirichlet boundary condition.

To solve this problem numerically, we first formulate it in a weak sense. We seek a solution $u$ in an appropriate
function space, $H^1(\Omega)$, which consists of functions that are square-integrable
and whose first derivatives are also square-integrable.  Multiplying the equation by a test function $v \in
    H^1(\Omega)$
and integrating  over $\Omega$, we obtain:
\[
    \int_\Omega \nabla u \cdot \nabla v \, dx - \int_{\Gamma} (\nabla u \cdot \mathbf{n}) \, v \, ds =
    \int_\Omega f \, v \, dx.
\]

We next introduce a triangulation \( \mathcal{T}_h \) consisting of quadrilateral (2D) or hexahedral (3D) elements of
size $h$, and define a finite element space $\mathbb{V}_h\subset H^1(\Omega)$ using Lagrange polynomial
elements of degree \( p \). In classical finite element methods, the mesh conforms to the boundary (unlike the
background mesh approach illustrated in Figure~\ref{fig:sbm_setup}), and test
functions $v$ typically vanish on $\Gamma $ to strongly enforce the Dirichlet condition $u=g$.
When function spaces do not necessarily satisfy essential boundary conditions strongly, boundary conditions
must be enforced weakly through additional integral terms on $\Gamma$.

Nitsche's method provides a way to weakly impose Dirichlet boundary conditions within a variational formulation without
requiring the function space to satisfy the boundary conditions. It modifies the bilinear form by adding terms on the
boundary $\Gamma$. The standard Nitsche formulation for the Poisson problem with Dirichlet boundary
conditions $u=g$ on $\Gamma$ seeks $u_h \in V_h$ such that for all $v_h \in \mathbb{V}_h$:
\begin{equation*}
    \begin{split}
        \int_\Omega \nabla u_h \cdot \nabla v_h \, dx
        \;\; - \int_{\Gamma} (\nabla u_h \cdot \mathbf{n}) v_h \, ds
        - &\alpha\int_{\Gamma} (\nabla v_h \cdot \mathbf{n}) u_h \, ds
        + \int_{\Gamma} \sigma u_h v_h \, ds = \\
        & = \int_\Omega f v_h \, dx
        - \int_{\Gamma} (\nabla v_h \cdot \mathbf{n}) g \, ds
        + \int_{\Gamma} \sigma_\Gamma g v_h \, ds.
    \end{split}
\end{equation*}
Here, $\sigma_\Gamma$ is a penalty parameter, typically chosen as $\sigma_\Gamma = \mathcal{O}(p^2 h^{-1})$ for mesh
size $h$, and
$\mathbf{n}$ is the outward unit normal vector to $\Gamma$.  The choice of parameter $\alpha=1$ leads to a symmetric
formulation, while $\alpha=-1$ results in a non-symmetric formulation in which the penalty term can be
skipped~\cite{baumann1999discontinuous}

\begin{figure}[!ht]
    \centering
\begin{tikzpicture}[scale=1.4]
    % Draw the Cartesian grid
    \draw[step=1cm, gray, very thin] (0.8,-0.2) grid (5.2,2.2);

    % Highlight cells inside the domain (interior cells)
    \foreach \x/\y in {2/1, 1/0, 1/1, 2/0, 3/0, 4/0}
        {
            \fill[green, opacity=0.5] (\x,\y) rectangle (\x+1,\y+1);
        }
    \foreach \x/\y in {1/0 }{
            \fill[green, opacity=0.5] (\x,\y) rectangle (\x-0.2,\y+2);
        }

    \fill[green, opacity=0.5] (0.8,-0.2) rectangle (5.2,0);

    \draw[line width=2pt, red, domain=51:90, samples=100, variable=\t]
    plot ({0.8 + 7.0*cos(\t)}, {-4.4 + 7.0*sin(\t)});
    \node at (2.6,2.5) {$\Gamma$};

    % Fill the area below the red line with light green to indicate the domain Omega
    \fill[green!30, opacity=0.4] plot[domain=51:90, samples=100, variable=\t] ({0.8 + 7.0*cos(\t)}, {-4.4 +
            7.0*sin(\t)}) -- (0.8,0) -- (5.2,0) -- cycle;

    % Define the angle parameter
    \pgfmathsetmacro{\myangle}{66}
    % Draw a vector from the surrogate boundary to the true boundary
    \draw[<-, thick] ({0.8 + 6.9*cos(\myangle) + 0.3}, {-4.4 + 6.9*sin(\myangle)}) -- ({0.8 + 5.9*cos(\myangle)
            +
            0.3}, {-4.4 +
            5.9*sin(\myangle)});
    % Position the label relative to the vector angle
    \pgfmathsetmacro{\labelangle}{\myangle - 6} % Adjust angle slightly for label placement
    \node at ({0.8 + 6.9*cos(\labelangle) - 0.6}, {-4.4 + 6.9*sin(\labelangle)+0.1}) {\large $\mathbf{d}$};

    \draw[->, thick] ({0.8 + 6.8*cos(\myangle-5) + 0.3}, {-4.4 + 6.8*sin(\myangle-5)}) -- ({0.8 +
            7.6*cos(\myangle-5)
            +
            0.3}, {-4.4 +
            7.6*sin(\myangle-5)});
    \node at ({0.8 + 7.5*cos(\labelangle-5) - 0.54}, {-4.4 + 7.5*sin(\labelangle-5)+0.36}) {\large $\mathbf{n}$};

    % Draw the vector d vertically from the surrogate boundary to the true boundary
    \draw[<-, thick]  ({0.8 + 5.9*cos(\myangle) +
            0.3}, {-4.4 + 6.9*sin(\myangle)}) -- ({0.8 + 5.9*cos(\myangle) +
            0.3}, {-4.4 +
            5.9*sin(\myangle)});
    \node at (3.4, 1.5) {\large $\tilde{\mathbf{n}}$}; % Place label near the normal vector

    % Draw the surrogate boundary (only the top boundary of the interior cells)
    \draw[line width=2pt, blue] (0.8,2) -- (3,2) -- (3,1) -- (4,1) -- (5,1)-- (5,0)-- (5.2,0.) ;

    \draw[line width=2pt, blue, dashed] (3,2) -- (4,2) -- (4,2) -- (4,1) ;
    \node at (1.7,2.2) {$\tilde{\Gamma}$};

    \node at (2.5,0.7) {$\tilde{\Omega}$};

    \node at (1.2,2.3) {${\Omega}$};

\end{tikzpicture}

    \caption{Schematic illustrating the background mesh, interior cells (green), the surrogate boundary
        $\tilde{\Gamma}$ (thick blue line) along the upper boundary of the interior cells, and the true boundary
        $\partial\Omega$ (red).}
    \label{fig:sbm_setup}
\end{figure}

\subsection{Shifted Boundary Method}
\mwerror{This section needs to be altered, for not it is a copy-paste from the previous paper.}
In many applications, the domain $\Omega$ may have a complex geometry, making the generation of body-fitted meshes
challenging. Non-body-fitted (unfitted) methods address this challenge by employing a background mesh $\mathcal{T}_h$
that
does not conform to the boundary of $\Omega$. The domain $\Omega$ is embedded within this background mesh, and cells of
$\mathcal{T}_h$ are classified as active based on their intersection with $\Omega$. The computational domain
$\tilde{\Omega}$ is defined as the union of these active cells. In the original SBM formulation~\cite{main2018shifted},
only cells strictly contained in $\Omega$ were included, leading to the surrogate domain boundary
depicted by the solid blue line in Figure~\ref{fig:sbm_setup}. Consequently, the surrogate domain $\tilde{\Omega}$ is
generally a subset of $\Omega$, and its boundary $\tilde{\Gamma}$ does not coincide with the true boundary $\Gamma$.

Later extensions include intersected cells~\cite{yang2024optimal} with a volume fraction
inside $\Omega$ greater than a threshold $\lambda\in [0,1]$. In Figure
\ref{fig:sbm_setup}, this
corresponds to including additional cells as indicated by the dashed blue
line. This approach reduces the distance between true and surrogate boundaries, but can lead to ill-conditioning.

To impose boundary conditions on $\tilde{\Gamma}$, the Dirichlet condition prescribed on the true boundary $\Gamma$ is
extrapolated to the surrogate boundary. This is typically accomplished by projecting points from $\tilde{\Gamma}$ onto
the true boundary $\Gamma$ along a suitable direction (e.g., the outward normal to $\tilde{\Gamma}$) and using a Taylor
expansion to approximate the boundary values. This is accomplished via an \emph{extension operator} $\mathcal{E}$,
which maps functions defined on $\tilde{\Omega}$ to the surrogate boundary $\tilde{\Gamma}$.

We assume that the Dirichlet boundary condition $g$ is given as a restriction of a function $u^\star$ defined on the
entire domain $\Omega$ to the boundary $\Gamma$. While the choice of this function $u^\star$ (as an extension of $g$
from $\Gamma$ into $\Omega$) is not unique, we take $u^\star$ to be equal to the solution $u$ in the
surrogate domain $\tilde{\Omega}$. For each point $\tilde{\mathbf{x}} \in \tilde{\Gamma}$, let $\mathbf{x} \in \Gamma$
be its closest point projection onto the true boundary, and let $\mathbf{d} = \mathbf{x} - \tilde{\mathbf{x}}$ be
the shift vector. The function $u$ is extended from $\Gamma$ to $\tilde{\Gamma}$ using a
Taylor expansion:
\[
    \mathcal{E}u^\star(\tilde{\mathbf{x}}) = u^\star(\mathbf{x}) - \mathbf{d} \cdot \nabla u (\tilde{\mathbf{x}}) +
    \cdots
\]
where $\mathcal{E}u^\star$ denotes the extrapolated boundary condition. The function $u^\star$ is assumed to be smooth
in a neighborhood of $\tilde{\Gamma}$, which allows for the Taylor expansion to be valid. By substituting it into the
weak formulation, we obtain a variational formulation for the shifted boundary problem with the extrapolated boundary
condition on
$\tilde{\Gamma}$ enforced in a Nitsche-like manner.

The SBM weak formulation seeks $u_h \in V_h$ such that for all $v_h
    \in V_h$,
\begin{equation*}
    \begin{split}
        \int_{\tilde{\Omega}} \nabla u_h \cdot \nabla v_h \, dx
        - \int_{\tilde{\Gamma}} (\nabla u_h \cdot \tilde{\mathbf{n}}) v_h \, ds
        - \alpha \int_{\tilde{\Gamma}} (\nabla v_h \cdot \tilde{\mathbf{n}}) \; \mathcal{E}u_h \, ds
        + \int_{\tilde{\Gamma}} \sigma_\Gamma \; \mathcal{E} u_h  \; v_h \, ds = \\
        = \int_{\tilde{\Omega}} f v_h \, dx
        - \int_{\tilde{\Gamma}} (\nabla v_h \cdot \tilde{\mathbf{n}}) g^E \, ds
        + \int_{\tilde{\Gamma}} \sigma_\Gamma g v_h \, ds.
    \end{split}
\end{equation*}
where $\tilde{\mathbf{n}}$ is the outward normal to $\tilde{\Gamma}$ and $\sigma_\Gamma$ is a penalty parameter.
We notice that our discrete solution $u_h$ is a piecewise polynomial function defined on the
background mesh $\mathcal{T}_h$, hence its Taylor expansion can be computed directly by evaluating the
function values at the points on the true boundary $\Gamma$. This allows us to avoid computation of higher-order
derivatives of $u_h$. Note that the resulting form is no longer symmetric due to the presence of the extrapolated
boundary condition. We will refer to formulation with $\alpha=1$ as the \emph{quasi-symmetric} formulation.

The choice of the stabilization term, particularly the parameter $\alpha$, significantly influences the spectral
properties of the resulting system matrix. As detailed in our previous work on a Discontinuous Galerkin (DG) based SBM
multigrid preconditioner~\cite{wichrowski2025geometric}, this can be observed even in a simple 1D single-cell problem.
For instance, a penalty-free formulation ($\sigma_\Gamma=0$, $\alpha=-1$) can lead to complex eigenvalues, especially
when the true boundary lies inside the surrogate domain.

In the context of DG methods, it was found that the choice of stabilization did not have a dramatic impact on the
overall multigrid performance, as the cell-based nature of the DG smoother effectively handled local
issues~\cite{wichrowski2025geometric}. However, in the DG formulation the cellwise smoother experienced singifficant
challenges for $p=3$. For continuous finite elements, where our smoother operate on
larger, overlapping patches of elements (i.e. vertex or edge patches), the choice of stabilization becomes far more
critical. The non-standard terms introduced by SBM are not well-contained within these patches, and an inappropriate
stabilization can introduce spectral properties that standard smoothers cannot handle effectively, a problem that is
exacerbated for higher-order elements.

\section{Multigrid Preconditioner}
\label{sec:mg_preconditioner}
The DG-SBM formulation described above leads to a large, sparse linear system. Due to
the lack of symmetry and possible indefiniteness of the resulting matrix, we employ a~Krylov subspace method,
specifically GMRES, for its solution.  The convergence of GMRES is sensitive to the condition number of the system
matrix that grows with both the number of elements in the mesh and the polynomial degree.
Multigrid methods~\cite{hackbusch1985multi} rely on a hierarchy of meshes with varying resolution levels. Our Cartesian
background mesh
naturally facilitates the construction of such a nested hierarchy. We define a sequence of
meshes $\{\mathcal{T}_\ell\}_{\ell=0}^L$, where $\ell$ denotes the level and $L$ is the finest level. For each level,
we identify active cells that have at least a fraction $\lambda$ of their measure (area in 2D, volume in 3D) inside the
domain.
On each mesh $\mathcal{T}_\ell$, we define a finite element space $\DGell$ consisting of piecewise polynomials of~
degree $p_\ell$.  Since the  set of active cells is determined by their intersection with the domain boundary, these
sets are not necessarily nested between levels, which can affect the efficiency of standard inter-grid transfer
operators. While our construction ensures a nested hierarchy  of finite element
spaces, the lack of a strict geometric relationship between active cells at different levels may reduce the
effectiveness of transfer operators. We mitigate this potential efficiency loss by using only linear elements on
geometrically coarser meshes and gradually increasing the polynomial degree only on the finest mesh. Hence the
resulting total  number of levels is $L+p$.

With this hierarchy of meshes and spaces established, the multigrid method is defined through two crucial
components: the smoothers, which reduce high-frequency error components on each level, and the transfer operators,
which move information between different resolution levels. These components, together with a coarse-grid solver, form
the complete multigrid algorithm. For preconditioning, we use a single multigrid V-cycle.

\subsection{Smoother}
\label{sec:smoother}
We first consider Richardson iterations with preconditioners $P$. Given the current
approximation $u_h$ and the right-hand side $f_h$, a~smoothing step updates
the approximation:
\[
    u_h^{k+1} = u_h^k + P(f_h - A_h u_h^k).
\]
We decompose the space $\DGell$ into a sum of subspaces $\mathbb{V}_i$, i.e., $\DGell = \sum_{i=1}^N
    \mathbb{V}_i$. Let $A_i$ be the restriction of $A$ to the subspace $\mathbb{V}_i$. Then, the additive subspace
correction
preconditioner is defined as:
\[
    P = \omega \sum_{i=1}^N A_i^{-1}
\]
where $\omega$ is a relaxation parameter. Alternatively, the successive subspace correction method is a subspace
correction method where the subspaces $\mathbb{V}_i$ are visited in a sequential manner. The procedure can be performed
by applying one Richardson iteration with the preconditioner $P_i$ for each subspace $\mathbb{V}_i$.
Since in each step only one subspace is corrected, the residual only changes locally and an efficient implementation is
possible. The preconditioner updates the solution by sequentially applying corrections for each subspace. For each
subspace $\mathbb{V}_i$, a correction is computed and applied:
\[
    u_h \leftarrow u_h + R_i^T A_i^{-1} R_i (f_h - A_h u_h).
\]
This process is repeated for all subspaces $i=1, \ldots, N$.
If the subspaces are chosen as one-dimensional spaces spanned by the basis functions, then the preconditioner is
equivalent to either Jacobi (additive) or Gauss-Seidel (successive).

\mwerror*{rewrite that part}{
    In the context of DG, a natural choice for the subspaces $\mathbb{V}_i$ is the set of functions that are non-zero
    only
    on a single cell. This leads to a cell-wise smoother, where we correct the solution on each cell independently.
    This
    can be interpreted as a block Gauss-Seidel smoother, where each block corresponds to the degrees of freedom
    associated
    with a single cell.}

Viewing the smoother as a subspace correction method reveals important insights into the convergence properties. The
correction computed within a subspace $\mathbb{V}_i$ implicitly satisfies certain boundary conditions on the boundary
of that subspace. This is particularly important in the context of the SBM, where the boundary conditions are shifted.
A careless choice of subspace can lead to overconstraining the correction, preventing certain error components from
being effectively smoothed. For instance, Jacobi and Gauss-Seidel smoothers for continuous finite element methods,
which
correspond to vertex-patch smoothing with zero correction at patch boundaries, conflict with the shifted boundary
conditions and are not effective, as confirmed in our preliminary numerical experiments. In contrast, the cell-wise
smoother aligns well with the DG method's element-based structure and takes into account the local nature of the SBM
boundary conditions, correcting the solution on each cell independently, as illustrated in Figure
\ref{fig:single_patch_boundaries}.

\begin{figure}[!ht]
    \centering
\begin{tikzpicture}[scale=1.4]
    % Draw the Cartesian grid
    \draw[step=1cm, gray, very thin] (0.8,-0.2) grid (5.2,2.2);

    % Highlight cells inside the domain (interior cells)
    \foreach \x/\y in {2/1, 1/0, 1/1, 2/0, 3/0, 4/0}
        {
            \fill[green, opacity=0.5] (\x,\y) rectangle (\x+1,\y+1);
        }
    \foreach \x/\y in {1/0 }{
            \fill[green, opacity=0.5] (\x,\y) rectangle (\x-0.2,\y+2);
        }

    \fill[green, opacity=0.5] (0.8,-0.2) rectangle (5.2,0);

    \draw[line width=2pt, red, domain=51:90, samples=100, variable=\t]
    plot ({0.8 + 7.0*cos(\t)}, {-4.4 + 7.0*sin(\t)});
    \node at (2.6,2.5) {$\Gamma$};

    % Fill the area below the red line with light green to indicate the domain Omega
    \fill[green!30, opacity=0.4] plot[domain=51:90, samples=100, variable=\t] ({0.8 + 7.0*cos(\t)}, {-4.4 +
            7.0*sin(\t)}) -- (0.8,0) -- (5.2,0) -- cycle;

    % Draw the surrogate boundary (only the top boundary of the interior cells)
    \draw[line width=1pt, blue] (0.8,2) -- (3,2) -- (3,1) -- (4,1) -- (5,1)-- (5,0)-- (5.2,0.) ;

    % Draw the surrogate boundary in the single patch case
    \draw[line width=3pt, blue] (2,2) -- (3,2) -- (3,1) -- (4,1);

    % Draw the interior boundary 
    \draw[line width=2pt, gray] (2,2) -- (2,0) -- (4,0) -- (4,1);

    \node at (1.5,0.7) {$\tilde{\Omega}$};

    \node at (1.5,2.3) {${\Omega}$};

\end{tikzpicture}

    \caption{
    Illustration of a patch. Interior faces ($\mathcal{F}_h^{\text{int}}$, green) connect to neighboring active cells.
    For cells adjacent to the domain boundary, faces can be part of the surrogate boundary $\tilde{\Gamma}$ (blue),
    which
    is shifted from
    the
    true boundary $\Gamma$ (red). The vertex-patch smoother solves should solve the  local problem on patch $P$ ...
    }
    \label{fig:single_patch_boundaries}
\end{figure}

\subsection{Full-Residual Shy Patches}

\subsection{Transfer Operators}
\label{sec:transfer_operators}
\mwerror{shorten and rewrite that part}
We use the standard $L^2$ projection operators provided by \texttt{deal.II} for the prolongation ($I_{\ell-1}^\ell$)
and restriction ($I_\ell^{\ell-1}$) operators, transferring data between coarser and finer grid levels. The~restriction
operator is the adjoint of the prolongation operator.

While these operators are standard in the context of multigrid methods, we note that the lack of a strict geometric
relationship between active cells at different levels may reduce the effectiveness of transfer operators. We tried to
reduce this potential efficiency loss by implementing a more sophisticated prolongation that would extrapolate
the values that are not defined on the coarser mesh, and its adjoint as a restriction. However, this approach did not
yield any significant improvements in the convergence rates. We therefore stick to the standard prolongation and
restriction operators.

The choice of the threshold parameter $\lambda$ directly influences the set of active cells, and thus the computational
domain $\tilde{\Omega}_\ell$, on each multigrid level $\ell$. A lower $\lambda$ value generally results in more cells
being classified as active, potentially leading to a surrogate domain that more closely envelops the true domain.
However, the sets of active cells across different multigrid levels, $\tilde{\Omega}_\ell$ and
$\tilde{\Omega}_{\ell-1}$, are not guaranteed to be nested. This lack of geometric nestedness can pose a challenge for
standard transfer operators. For instance, a coarse cell might be active while some of its children on the finer level
are not, or vice-versa. This mismatch can reduce the efficiency of inter-grid transfers, as information might be
prolonged to or restricted from regions that are not consistently part of the active domain across levels. While our
experiments with more sophisticated extrapolation-based transfer operators did not show significant gains, the
interplay between $\lambda$, the resulting active cell configurations, and the transfer operator effectiveness remains
an important consideration for multigrid performance in SBM.

\section{Implementation details}
\label{sec:implementation}
\mwerror{refer to the previous paper, say that we follow it. Shorten.}
The numerical implementation of the DG-SBM multigrid solver is based on the open-source finite element library
\texttt{deal.II}~\cite{dealii2024}. It provides a comprehensive framework for the implementation of~finite
element methods, including mesh handling, finite element spaces, assembly of linear systems, and interfaces to various
linear algebra solvers and preconditioners. Its flexibility and extensive capabilities make it well-suited for
developing complex methods like the one presented here, which involves unfitted meshes and multigrid techniques.

The spatial discretization of the Poisson equation using the Discontinuous Galerkin method follows the standard
framework available within \texttt{deal.II}. This involves defining the DG finite element space on the background
Cartesian mesh, assembling the element-wise contributions to the global stiffness matrix and right-hand side vector,
and handling the jump and average terms across interior faces as described in Section 2.2. The library's support for
various DG formulations simplifies the implementation of the bilinear form $a_h(u_h, v_h)$.

\subsection{Background mesh and geometry handling}
When implementing SBM on a non-body-fitted mesh, we need to handle cells intersected by the true
boundary $\Gamma$. We use a level set function $\phi(\mathbf{x})$ to implicitly define the domain $\Omega =
    \{\mathbf{x} \mid \phi(\mathbf{x}) < 0\}$, with $\Gamma = \{\mathbf{x} \mid \phi(\mathbf{x}) = 0\}$. Cells of the
background mesh are classified based on their intersection with the zero level set: cells entirely inside $\Omega$
(interior), cells entirely outside $\Omega$ (exterior), and cells intersected by $\Gamma$.
Then, for the intersected cells the fraction of the cell volume inside $\Omega$ is computed, and the cell is
classified as active if this fraction is greater than a threshold $\lambda$. This is handled using non-matching
quadrature rules implemented in \texttt{deal.II}, which are based on the techniques described in~\cite{saye2015high}.

In our approach, degrees of freedom are formally assigned to all cells of the background mesh,
including those that are classified  as non-active. However, for non-active cells, the local contribution to the global
system matrix is set to the identity matrix, and the corresponding entries in the right-hand side are set to zero. This
effectively decouples the
degrees of freedom in non-active cells from the system, ensuring that the solution is only computed within the
computational domain $\tilde{\Omega}$ while simplifying the global matrix assembly process by maintaining a uniform
structure of degrees of freedom (DoF) across all cells.

\subsection{Processing surrogate boundary and matrix assembly}
\label{sec:sbm_assembly}
The SBM requires computing the closest point projection from points on the surrogate boundary $\tilde{\Gamma}$ to the
true boundary $\Gamma$. For the specific test cases, where it is possible, this projection can be computed exactly.
For~
general geometries, finding the closest point involves solving a local nonlinear optimization problem for each
point on $\tilde{\Gamma}$.
The problem is to find the point on
the true boundary that minimizes the distance to a given point on the surrogate boundary, since the true boundary is
given by the zero level set of the level set function $\phi(\mathbf{x})$, this can be formulated as:
\[
    \min_{\mathbf{x} \in \Gamma} \|\mathbf{x} - \tilde{\mathbf{x}}\|^2
    \quad \text{s.t.} \quad \phi(\mathbf{x}) = 0.
\]
We introduce a Lagrange multiplier $\mu$ to enforce the constraint $\phi(\mathbf{x}) = 0$. The Lagrangian function
is:
\[
    \mathcal{L}(\mathbf{x}, \mu) = \|\mathbf{x} - \tilde{\mathbf{x}}\|^2 + \mu \phi(\mathbf{x}).
\]
The necessary conditions for optimality are given by the stationary points of the Lagrangian, which satisfy the
following system of equations:
\begin{align*}
    \nabla_{\mathbf{x}} \mathcal{L}(\mathbf{x}, \mu)           & = 2(\mathbf{x} - \tilde{\mathbf{x}}) + \mu
    \nabla\phi(\mathbf{x}) = 0,
    \\
    \frac{\partial \mathcal{L}}{\partial \mu}(\mathbf{x}, \mu) & = \phi(\mathbf{x}) = 0.
\end{align*}
Solving this system yields the closest point $\mathbf{x}$ on the true boundary $\Gamma$ to the point
$\tilde{\mathbf{x}}$ on the surrogate boundary $\tilde{\Gamma}$. The nonlinear system is solved
iteratively using a Newton-Raphson method. As the Newton-Raphson method requires the evaluation of the Jacobian matrix,
we need to compute the second derivatives of the level set function $\phi(\mathbf{x})$. To ensure sufficient smoothness
for derivative calculations in the closest point search, the level set function is represented using finite elements of
degree $2$ for $p=1$
and degree $p$ for $p>1$.

The search for the closest point is performed only in the interior of the cells adjacent to the surrogate
boundary. While this does not guarantee that the closest point is found inside the cell, it is expected that
even if the closest point is outside the cell, the extrapolation will still yield a good approximation of the boundary
condition. We will verify this assumption in the numerical results.

Finally, the extrapolation of the function values from the true boundary to the
surrogate boundary required for the matrix assembly process is accomplished by evaluating the values of the basis
functions at the points on the surrogate boundary. The assembly of the cell contributions and interior faces to the
system
matrix and right-hand side vector is performed using the standard finite element assembly process provided by
\texttt{deal.II}.

\subsection{Multigrid structures}

For the multigrid hierarchy, we leverage \texttt{deal.II}'s built-in capabilities for handling nested meshes and
defining finite element spaces on each level. The standard $L^2$ projection operators provided by \texttt{deal.II} are
used for the prolongation ($I_{\ell-1}^\ell$) and restriction ($I_\ell^{\ell-1}$) operators, transferring data between
coarser and finer grid levels.

The assembly of the system matrix $A_\ell$ on each level $\ell$ of the multigrid hierarchy follows a procedure
analogous to that on the finest level. The level set function defining the domain geometry is interpolated onto the
mesh $\mathcal{T}_\ell$. Based on this interpolated level set and the chosen threshold $\lambda$, active cells for
level $\ell$ are identified. The DG-SBM bilinear form is then used to assemble the local contributions to $A_\ell$ only
for these active cells. For cells deemed non-active on level $\ell$, their corresponding entries in $A_\ell$ are set to
effectively decouple them from the system by contributing an identity block to the matrix and zero to the right-hand
side.

The smoother component of the multigrid V-cycle is implemented as a cell-wise Symmetric Successive Over-Relaxation
(SSOR) method.	The SSOR iteration proceeds by visiting each cell sequentially and solving a small local linear system
corresponding to the block of degrees of freedom within that cell, using the most recently updated values from
previously processed cells.
This choice is motivated by its relative ease of tuning compared to other smoothers. For instance, a Block Jacobi
smoother typically requires a carefully chosen relaxation parameter. In contrast, SSOR-based smoothers
often perform well with a fixed relaxation parameter, frequently $\omega=1$, reducing the need for problem-specific
tuning. Unless stated otherwise, we use $\omega=1$.

The current implementation focuses on testing the fundamental concepts and effectiveness of the multigrid
preconditioner for the DG-SBM method. Therefore, it is based on serial computations using sparse matrices. While
\texttt{deal.II} supports parallelization, the initial goal was to validate the multigrid algorithm's convergence
properties and identify potential challenges in the DG-SBM context before scaling up to larger, parallel problems.

\section{Numerical results}
\label{sec:numerical_results}
In this section, we present numerical results to evaluate the performance of the proposed DG-SBM multigrid
preconditioner for the Poisson equation. We investigate its effectiveness in terms of convergence rates, iteration~
counts under mesh refinement ($h$-refinement) and polynomial degree increase ($p$-refinement). To better illustrate the
multigrid preconditioner's performance characteristics and enable meaningful comparison of iteration counts, we solve
the system to high accuracy with a tolerance of $10^{-12}$ for the relative residual reduction. This stringent
tolerance results in higher iteration counts, making the differences in solver performance more apparent and
facilitating clearer assessment of the multigrid effectiveness. Solver failure is defined as exceeding 100 GMRES
iterations without reaching this tolerance.
The initial background mesh is a square (or cube in 3D) domain covering the range $[-1.01, 1.01]$ in each dimension,
subdivided into $4$ cells in each coordinate direction. We use $\sigma_\Gamma = 5$ and $\sigma_F = 1$ for the penalty
parameters in the SBM
formulation. The interior penalty $\sigma_F=1$ is chosen as a minimal value to ensure coercivity while minimizing its
impact on solver convergence, whereas the boundary penalty $\sigma_\Gamma=5$ is selected to be sufficiently large to
ensure stability of the Nitsche coupling on the surrogate boundary.

The majority of our tests are conducted on a unit sphere (a unit disk in 2D, depicted on the left panel of
Fig.~\ref{fig:shift_distribution}). While geometrically simple, the unit
sphere serves as an insightful benchmark. Any sufficiently smooth ($C^1$) complex boundary, when viewed at a fine
enough mesh resolution, locally resembles a flat plane. The unit sphere, due to its uniform curvature, presents a
comprehensive range of intersection angles between the true boundary and the background mesh cells. Furthermore, the
boundary intersects cells at various locations relative to cell centers and faces, leading to a diverse distribution of
shift vector magnitudes and directions. This includes scenarios where the shift vector points from the surrogate
boundary towards the interior of the true domain (which we denote as negative shifts if they oppose the outward normal
of the surrogate boundary). The right panel of Figure~\ref{fig:shift_distribution}
depicts the minimum and maximum shift magnitudes observed across the surrogate boundary for a typical
discretization. The shift magnitudes are normalized by the cell size $h$; for a unit cell, the theoretically largest
possible shift magnitude in 2D is $\sqrt{2}h$, occurring if the surrogate boundary point is at a cell corner and the
true boundary passes through the diagonally opposite corner.

The primary metric for evaluating the multigrid preconditioner is the number of GMRES iterations required to reduce the
initial residual by a factor of $10^{-12}$. If the threshold is not met within 100 iterations, the solver is considered
to have failed. All experiments are conducted using the implementation described in the previous section. As our solver
depends on the value of the threshold parameter $\lambda$, we will first explore the solver performance on 2D problems
as they are less computationally intensive and then show some results for 3D problems with tuned $\lambda$.

\section{Conclusion}
\label{sec:conclusion}

\section*{Acknowledgments}
% Optional: Acknowledge funding sources, helpful discussions, etc.
The author declares the use of language models (ChatGPT, Gemini, and Claude) to improve the clarity and readability of
the manuscript. All scientific content and technical claims are solely the responsibility of the author.

The author would like to thank Guglielmo Scovazzi and Guido Kanschat for insightful discussions, and Luca Heltai for
encouraging the exploration of this method.

\FloatBarrier
\bibliographystyle{siam} % Or any other style
\bibliography{literature} % Assuming a references.bib file will be created 
\end{document}